\documentclass[11pt]{article}

\usepackage[a4paper,margin=1in]{geometry}
\usepackage[T1]{fontenc}
\usepackage{lmodern}
\usepackage{microtype}
\setlength{\emergencystretch}{1em}
\raggedbottom
\usepackage{amsmath,amssymb,amsthm,mathtools}
\usepackage{booktabs}
\usepackage{graphicx}
\usepackage{xcolor}
\usepackage{hyperref}
\usepackage[nameinlink,noabbrev]{cleveref}

\hypersetup{
  colorlinks=true,
  linkcolor=blue!55!black,
  citecolor=blue!55!black,
  urlcolor=blue!55!black,
  pdftitle={Half-Period Involutions and Exact Cancellation in Inert Lucas Sequences},
  pdfauthor={Majid Ghandali}
}

\newcommand{\ord}{\operatorname{ord}}
\newcommand{\Norm}{\operatorname{Norm}}
\newcommand{\Fp}{\mathbb{F}_p}
\newcommand{\Fpq}{\mathbb{F}_{p^2}}
\newcommand{\Z}{\mathbb{Z}}
\newcommand{\chiP}{\chi_p}

\theoremstyle{plain}
\newtheorem{theorem}{Theorem}[section]
\newtheorem{proposition}[theorem]{Proposition}
\newtheorem{lemma}[theorem]{Lemma}
\newtheorem{corollary}[theorem]{Corollary}

\theoremstyle{definition}
\newtheorem{definition}[theorem]{Definition}
\newtheorem{example}[theorem]{Example}
\newtheorem{remark}[theorem]{Remark}

\title{Half-Period Involutions and Exact Cancellation in Inert Lucas Sequences}

\author{
  Majid Ghandali\\
  Independent Researcher, Tehran, Iran\\
  \href{mailto:majid.ghandali@gmail.com}{\texttt{majid.ghandali@gmail.com}}\\
 ORCID:
\href{https://orcid.org/0009-0001-1097-1770}
{0009-0001-1097-1770}
}
\date{}

\begin{document}

\maketitle

\begin{abstract}
We organize two elementary and exact cancellation mechanisms for
quadratic-character sums attached to Lucas sequences modulo an inert odd
prime. The first is translation by a matrix half-period, yielding the
anti-periodicity relation $U_{n+T/2}=-U_n$. The second is rank-block scalar
propagation, equivalently a geometric character factorization, in which the
multiplier $\Lambda^k$ filters complete character sums. These mechanisms
arise from different structural viewpoints: the former is an odd-function
pairing under index translation, whereas the latter is multiplicative
filtering by the rank scalar. The determinant/norm relation links the
companion-matrix determinant to the field norm of a characteristic root and
forces $\ord(Q)\mid T$; it is kept as an independent observation and is not
itself a cancellation criterion. In the quadratic-character full-period
setting, the two criteria are nevertheless related: the half-period
quadratic criterion implies the rank-block scalar criterion, while scalar
cancellation may occur even when the half-period negation is invisible to
$\chiP$. The contribution is the explicit separation and comparison of
these elementary mechanisms within the inert finite-field framework.
\end{abstract}

\noindent\textbf{Keywords:} Lucas sequences; rank of apparition; matrix period; finite-field characters; quadratic characters; exact cancellation.

\noindent\textbf{MSC 2020:} 11B39, 11T24, 11L10.

\section{Introduction}

Let $(U_n)_{n\geq 0}$ be the Lucas sequence
\[
 U_0=0,\qquad U_1=1,\qquad
 U_{n+2}=P U_{n+1}-Q U_n,
 \qquad P,Q\in\Z.
\]
For an odd prime $p$, we study complete-period sums of the form
\[
 \sum_n \chiP(U_n),
\]
where $\chiP$ is the quadratic character of $\Fp$, extended by $\chiP(0)=0$.

The arithmetic of Lucas sequences goes back to Lucas and Lehmer
\cite{Lucas1878,Lehmer1930}; for historical and general accounts, see
\cite{Williams,Everest2003}. Periods modulo $m$, ranks of apparition or
entry points, and associated order phenomena for Fibonacci and Lucas-type
recurrences are treated in the classical works
\cite{Wall1960,Robinson1963,Vinson1963,Somer1980} and in later developments
\cite{Renault1996,Renault2013,FiebigMbirikaSpilker2025}. Divisibility and
rank/index questions for Lucas and Lehmer sequences form a broader related
literature \cite{BHV2001,Stewart2013,Sanna2022,Sanna2024}. The finite-field
background used below is standard \cite{LN}; for a splitting-field
perspective on Fibonacci periods modulo primes, see
\cite{GuptaRockstrohSu2012}. Character sums for recurrence sequences form
another related strand; see, for example, \cite{BlackburnShparlinski2006}.

The purpose of this concise note is structural. We separate two cancellation
arguments that are often used in different forms. One is a half-period
involution of the companion matrix. The other is rank-block scalar
propagation, or equivalently geometric character factorization, arising from
the scalar multiplier of each rank block. The determinant/norm relation is
kept separate as an independent observation about the matrix period. The
distinction matters, especially when $\chiP(-1)=1$ but $\chiP(\Lambda)=-1$.
Among recent contributions, Fiebig, Mbirika, and Spilker
\cite{FiebigMbirikaSpilker2025} develop a period--entry-point framework for
Lucas sequences modulo $m$, including structural relations among periods,
entry points, and associated orders. The present work considers a more
specialized inert-prime setting, where the companion matrix admits a
quadratic-extension interpretation. In this setting, the multiplier
associated with the rank of apparition provides an additional scalar
structure that leads to exact character-sum factorization. Thus, the present
approach complements the period--entry-point viewpoint by exploiting
multiplicative-order information available in the inert quadratic-extension
framework.

The contribution of this note is the precise separation of two logically distinct cancellation mechanisms for character sums in the inert Lucas setting. The first is an index-translation mechanism induced by the half-period involution $A^{T/2}=-I$. The second is a scalar rank-block mechanism induced by $A^\alpha=\Lambda I$, leading to a geometric character factor. We show that, for quadratic characters over a full matrix period, the former criterion implies the latter, whereas the converse may fail as a cancellation proof when $\chi_p(-1)=1$. This distinction is made explicit through the finite-field multiplier $\Lambda$ and through concrete examples.

\section{Setup}

Throughout, let $P,Q\in\Z$, let $p$ be an odd prime, and put
\[
 D=P^2-4Q.
\]
Although Lucas sequences are often studied under characteristic-zero non-degeneracy assumptions, the finite-field arguments of this note require only the local conditions
\[
p\nmid QD
\]
and inertness. We use the following standard terminology.

\begin{definition}
The prime $p$ is called \emph{inert} for the sequence if
\[
 \chiP(D)=-1;
\]
equivalently, the polynomial
\[
 f(X)=X^2-PX+Q\in\Fp[X]
\]
is irreducible.
\end{definition}

In the inert case,
\[
 \Fpq=\Fp(\sqrt D),
\]
and the two roots of $f$ are
\[
 \lambda=\frac{P+\sqrt D}{2},
 \qquad
 \mu=\frac{P-\sqrt D}{2}.
\]
They satisfy
\[
 \lambda+\mu=P,\qquad \lambda\mu=Q,
 \qquad \mu=\lambda^p,\qquad \lambda=\mu^p.
\]
The Binet formulas are
\begin{equation}
 U_n=\frac{\lambda^n-\mu^n}{\lambda-\mu},
 \qquad
 V_n=\lambda^n+\mu^n,
 \label{eq:binet}
\end{equation}
where
\[
 V_0=2,\qquad V_1=P,\qquad
 V_{n+2}=P V_{n+1}-Q V_n.
\]
Both sequences are regarded modulo $p$.

Let
\[
 A=
 \begin{pmatrix}
 P&-Q\\
 1&0
 \end{pmatrix}
 \in\operatorname{GL}_2(\Fp).
\]
For $n\geq1$,
\begin{equation}
 A^n=
 \begin{pmatrix}
 U_{n+1}&-Q U_n\\
 U_n&-Q U_{n-1}
 \end{pmatrix}.
 \label{eq:matrix-power}
\end{equation}
In particular, the lower-left entry of $A^n$ is $U_n$. This companion-matrix viewpoint is classical; for matrix and period/rank approaches to Fibonacci and related recurrences modulo $m$, see \cite{Robinson1963,Wall1960,Vinson1963,Renault2013}.

\section{Matrix Period}

Define
\[
 T=\ord_{\operatorname{GL}_2(\Fp)}(A).
\]
Since $A^T=I$, the sequences $(U_n)_{n\geq0}$ and $(V_n)_{n\geq0}$ are periodic modulo $p$, and $T$ is a common period for both sequences. The minimal periods of the individual sequences may divide $T$.

\begin{proposition}[Order of the eigenvalues]
\label{prop:eigenvalue-order}
In the inert case,
\[
 T=\ord_{\Fpq^\times}(\lambda)
  =\ord_{\Fpq^\times}(\mu).
\]
Consequently,
\[
 T\mid p^2-1.
\]
\end{proposition}

\begin{proof}
The matrix $A$ is diagonalizable over $\Fpq$, with eigenvalues $\lambda$ and $\mu$. Since $\mu=\lambda^p$ and the Frobenius map $x\mapsto x^p$ is an automorphism of $\Fpq^\times$, the two eigenvalues have the same multiplicative order. Hence the order of $A$ is their common order. This order divides
\[
 |\Fpq^\times|=p^2-1.
\]
\end{proof}

\begin{proposition}[Half-period criterion]
\label{prop:half-period-criterion}
For an integer $m$, one has
\[
 A^m=-I
\]
if and only if $T$ is even and
\[
 m\equiv \frac{T}{2}\pmod T.
\]
In particular, if $T$ is even, then
\[
 A^{T/2}=-I
 \qquad\text{and}\qquad
 U_{n+T/2}=-U_n,\quad
 V_{n+T/2}=-V_n
\]
for every nonnegative integer $n$.
\end{proposition}

\begin{proof}
By \cref{prop:eigenvalue-order}, $A^m=-I$ is equivalent to
\[
 \lambda^m=\mu^m=-1.
\]
Since $\lambda$ and $\mu$ both have order $T$, this occurs precisely when $T$ is even and
\[
 m\equiv T/2\pmod T.
\]
For $m=T/2$, the identity $A^{T/2}=-I$ and the matrix formula \eqref{eq:matrix-power} give the assertion for $U_n$. Alternatively, the Binet formulas \eqref{eq:binet} give both anti-periodicity relations directly.
\end{proof}

\begin{remark}[Determinant and norm]
\label{rem:determinant-norm}
Since
\[
 \det(A)=Q,
\]
the relation $A^T=I$ implies
\[
 Q^T=1\pmod p,
 \qquad
 \ord_{\Fp^\times}(Q)\mid T.
\]
The same relation has the field-theoretic interpretation
\[
 \Norm_{\Fpq/\Fp}(\lambda)=\lambda\lambda^p
 =\lambda\mu=Q.
\]
Thus the determinant of the companion matrix is the norm of either characteristic root. This is an independent observation and does not itself imply a half-period involution.
\end{remark}

\section{Half-Period Translation and Determinant/Norm Observation}

The half-period relation in \cref{prop:half-period-criterion} is a translation statement on indices. It should not be confused with the determinant or norm constraint in \cref{rem:determinant-norm}. The latter forces the order of $Q$ to divide the matrix period, but does not by itself imply that the period is even or that its midpoint acts as $-I$.

For a function $g:\Fp\to\mathbb{C}$ satisfying
\[
 g(-x)=-g(x),
\]
the half-period relation immediately yields exact cancellation on any index set invariant under the translation $n\mapsto n+T/2$.

\begin{proposition}[Half-period cancellation]
\label{prop:half-period-cancellation}
Assume that $T$ is even. Let
$I\subseteq\Z/T\Z$ be stable under
\[
 n\longmapsto n+\frac{T}{2}.
\]
If $g:\Fp\to\mathbb{C}$ is odd, then
\[
 \sum_{n\in I}g(U_n)=0
 \qquad\text{and}\qquad
 \sum_{n\in I}g(V_n)=0.
\]
\end{proposition}

\begin{proof}
The translation $n\mapsto n+T/2$ partitions $I$ into pairs. By \cref{prop:half-period-criterion},
\[
 g(U_{n+T/2})=g(-U_n)=-g(U_n),
\]
and the same argument applies to $V_n$. Summing over the pairs proves both identities.
\end{proof}

\begin{corollary}[Quadratic half-period cancellation]
\label{cor:quadratic-half-period}
If $T$ is even, $I$ is stable under $n\mapsto n+T/2$, and
\[
 \chiP(-1)=-1,
\]
then
\[
 \sum_{n\in I}\chiP(U_n)=0
 \qquad\text{and}\qquad
 \sum_{n\in I}\chiP(V_n)=0.
\]
\end{corollary}

\begin{proof}
The extension $\chiP(0)=0$ preserves the identity
\[
 \chiP(-x)=\chiP(-1)\chiP(x)=-\chiP(x)
\]
for every $x\in\Fp$. Apply \cref{prop:half-period-cancellation}.
\end{proof}

\section{Rank Blocks and Scalar Character Filtering}

Let
\[
 \alpha=\min\{n\geq1:U_n=0\pmod p\}
\]
be the rank of apparition of $p$, and define
\[
 \Lambda=U_{\alpha+1}\pmod p.
\]
In the inert setting the set above is nonempty, so $\alpha$ is well defined: \cref{lem:rank-characterization} below identifies $\alpha$ with the order of $\lambda/\mu\in\Fpq^\times$. (For ranks of apparition in broader Lucas-sequence settings, see \cite{Desmond1978,Somer1980,Sanna2022,Sanna2024}.)

\begin{lemma}[Rank characterization]
\label{lem:rank-characterization}
In the inert case,
\[
 U_n=0\pmod p
 \quad\Longleftrightarrow\quad
 \left(\frac{\lambda}{\mu}\right)^n=1.
\]
Consequently,
\[
 \alpha=\ord_{\Fpq^\times}\left(\frac{\lambda}{\mu}\right).
\]
\end{lemma}

\begin{proof}
Since $\lambda-\mu\neq0$ in $\Fpq$, the Binet formula gives
\[
 U_n=0
 \quad\Longleftrightarrow\quad
 \lambda^n-\mu^n=0
 \quad\Longleftrightarrow\quad
 \left(\frac{\lambda}{\mu}\right)^n=1.
\]
The assertion follows from the definition of $\alpha$.
\end{proof}

\begin{corollary}
\label{cor:alpha-divides-p-plus-1}
In the inert case, $\alpha\mid p+1$.
\end{corollary}

\begin{proof}
Since $\mu=\lambda^p$ and $\lambda\mu=Q$, one has $\lambda^{p+1}=\mu^{p+1}=Q$, hence
\[
 \left(\frac{\lambda}{\mu}\right)^{p+1}=1.
\]
Therefore $\ord_{\Fpq^\times}(\lambda/\mu)$ divides $p+1$, and \cref{lem:rank-characterization} gives $\alpha=\ord_{\Fpq^\times}(\lambda/\mu)$.
\end{proof}

\begin{proposition}[Rank multiplier and period decomposition]
\label{prop:rank-multiplier}
One has
\[
 A^\alpha=\Lambda I
 \qquad\text{and}\qquad
 \lambda^\alpha=\mu^\alpha=\Lambda.
\]
If
\[
 \omega=\ord_{\Fp^\times}(\Lambda),
\]
then
\begin{equation}
 T=\alpha\omega.
 \label{eq:period-decomposition}
\end{equation}
Moreover, for every $k\geq0$ and $1\leq j\leq\alpha$,
\[
 A^{k\alpha+j}=\Lambda^kA^j,
 \qquad
 U_{k\alpha+j}=\Lambda^kU_j,
 \qquad
 V_{k\alpha+j}=\Lambda^kV_j.
\]
\end{proposition}

\begin{proof}
Using $U_\alpha=0$ in \eqref{eq:matrix-power},
\[
 A^\alpha=
 \begin{pmatrix}
 U_{\alpha+1}&0\\
 0&-Q U_{\alpha-1}
 \end{pmatrix}.
\]
The recurrence at index $\alpha-1$ gives
\[
 U_{\alpha+1}=P U_\alpha-Q U_{\alpha-1}
             =-Q U_{\alpha-1},
\]
so $A^\alpha=\Lambda I$. Since $A^\alpha$ has eigenvalues $\lambda^\alpha$ and $\mu^\alpha$, and equals the scalar matrix $\Lambda I$, it follows that $\lambda^\alpha=\mu^\alpha=\Lambda$.

By \cref{lem:rank-characterization}, $U_n\equiv0\pmod p$ if and only if $(\lambda/\mu)^n=1$, and since $\alpha=\ord(\lambda/\mu)$, this holds precisely when $\alpha\mid n$. Now $A^n=I$ implies $U_n=0$ by the lower-left entry of \eqref{eq:matrix-power}, hence $\alpha\mid n$; writing $n=k\alpha$,
\[
 A^n=(A^\alpha)^k=\Lambda^k I,
\]
which equals $I$ exactly when $\omega\mid k$. Thus $A^n=I$ precisely when $\alpha\omega\mid n$, and \eqref{eq:period-decomposition} follows.

Finally,
\[
 A^{k\alpha+j}
 =(A^\alpha)^kA^j
 =\Lambda^kA^j,
\]
which gives the formula for $U_{k\alpha+j}$ by comparison of lower-left entries, and the formula for $V_{k\alpha+j}$ follows from the Binet formula \eqref{eq:binet}:
\[
 V_{k\alpha+j}
 =\lambda^{k\alpha+j}+\mu^{k\alpha+j}
 =\Lambda^k(\lambda^j+\mu^j)
 =\Lambda^kV_j.
\]
\end{proof}

\begin{remark}[Classical context]
\label{rem:classical-context}
The identities in \cref{prop:rank-multiplier} should be viewed as the inert-prime
specialization of classical period--rank--order relations for Fibonacci and
Lucas-type sequences; compare
\cite{Wall1960,Robinson1963,Vinson1963,Somer1980,Renault1996,Renault2013,FiebigMbirikaSpilker2025}.
We include the proof to fix the notation $\alpha,\Lambda,\omega$ and to make the
subsequent character-factorization argument self-contained. The divisibility
$\alpha\mid p+1$ in \cref{cor:alpha-divides-p-plus-1} is classical in the nonsplit
prime setting as well \cite{Lucas1878,Wall1960,GuptaRockstrohSu2012}.
Taking determinants in $A^\alpha=\Lambda I$ further gives
\[
\Lambda^2=Q^\alpha,
\]
which connects the rank multiplier with the determinant/norm relation of
\cref{rem:determinant-norm}.
In the present inert prime setting, the multiplier admits an explicit
finite-field interpretation through the characteristic root, and its
multiplicative order gives the scalar decomposition underlying the rank-block
character factorization developed below.
\end{remark}

\begin{proposition}[Full-period character factorization]
\label{prop:full-period-factorization}
Let $\psi$ be any multiplicative character of $\Fp^\times$, extended by $\psi(0)=0$. For $X=U$ or $X=V$,
\begin{equation}
 \sum_{n=1}^{T}\psi(X_n)
 =
 \left(\sum_{k=0}^{\omega-1}\psi(\Lambda)^k\right)
 \left(\sum_{j=1}^{\alpha}\psi(X_j)\right).
 \label{eq:full-period-factorization}
\end{equation}
In particular, if
\[
 \psi(\Lambda)\neq1,
\]
then
\[
 \sum_{n=1}^{T}\psi(U_n)=0
 \qquad\text{and}\qquad
 \sum_{n=1}^{T}\psi(V_n)=0.
\]
\end{proposition}

\begin{proof}
Partition the complete period into the blocks
\[
 n=k\alpha+j,
 \qquad
 0\leq k\leq\omega-1,\quad 1\leq j\leq\alpha.
\]
By \cref{prop:rank-multiplier},
\[
 X_{k\alpha+j}=\Lambda^kX_j.
\]
Therefore,
\[
 \psi(X_{k\alpha+j})
 =\psi(\Lambda)^k\psi(X_j).
\]
Summing first over $j$ and then over $k$ proves \eqref{eq:full-period-factorization}. If $\psi(\Lambda)\neq1$, the first factor is a geometric sum over a complete cycle of $\psi(\Lambda)$ and is zero.
\end{proof}

Thus the full-period character sum separates into an inner rank-block component and an outer scalar-orbit component. The resulting cancellation may occur entirely in the scalar-orbit direction, independently of any sign-reversing involution on the index set.

\begin{corollary}[Quadratic scalar cancellation]
\label{cor:quadratic-scalar}
If
\[
 \chiP(\Lambda)=-1,
\]
then $\omega$ is even and
\[
 \sum_{n=1}^{T}\chiP(U_n)=0
 \qquad\text{and}\qquad
 \sum_{n=1}^{T}\chiP(V_n)=0.
\]
\end{corollary}

\begin{proof}
The condition $\chiP(\Lambda)=-1$ implies that $\Lambda$ has even order (if $\omega$ were odd, then $\Lambda=(\Lambda^{(\omega+1)/2})^2$ would be a square), so $\omega$ and $T=\alpha\omega$ are even. The vanishing follows from \cref{prop:full-period-factorization}, since $\chiP(\Lambda)\neq1$.
\end{proof}

\begin{proposition}[Quadratic character of the rank multiplier]
\label{prop:chi-lambda}
In the inert setting:
\begin{enumerate}
 \item[(a)] if $\alpha$ is even, then $\chiP(\Lambda)=-1$;
 \item[(b)] if $p\equiv3\pmod4$, then $\chiP(\Lambda)=-1$ if and only if $T$ is even.
\end{enumerate}
In particular, whenever the half-period criterion of \cref{cor:quadratic-half-period} applies to a full period, the rank-block scalar criterion of \cref{cor:quadratic-scalar} applies as well; the converse need not hold (\cref{ex:scalar-cancellation}).
\end{proposition}

\begin{proof}
(a) Put $r=\lambda/\mu$, so that $\ord(r)=\alpha$ by \cref{lem:rank-characterization}, and recall that $\lambda^\alpha=\Lambda$ by \cref{prop:rank-multiplier}. By Euler's criterion, since $\Lambda\in\Fp^\times$,
\[
 \chiP(\Lambda)
 =\Lambda^{(p-1)/2}
 =\lambda^{\alpha(p-1)/2}
 =\bigl(\lambda^{p-1}\bigr)^{\alpha/2}
 =\bigl(\mu/\lambda\bigr)^{\alpha/2}
 =r^{-\alpha/2}.
\]
Since $\alpha$ is even and $\ord(r)=\alpha$, the element $r^{\alpha/2}$ has order $2$ in the cyclic group $\Fpq^\times$, hence equals $-1$; therefore $r^{-\alpha/2}=-1$ and $\chiP(\Lambda)=-1$.

(b) Suppose first that $T$ is even, so that $\alpha\omega$ is even. If $\alpha$ is even, then $\chiP(\Lambda)=-1$ by (a). If $\omega$ is even, then $\Lambda$ cannot be a square: otherwise $\omega=\ord(\Lambda)$ would divide the order $(p-1)/2$ of the subgroup of squares of $\Fp^\times$, which is odd when $p\equiv3\pmod4$, contradicting the evenness of $\omega$. Hence $\chiP(\Lambda)=-1$.

Conversely, suppose $\chiP(\Lambda)=-1$. If $\omega$ were odd, then $\Lambda=(\Lambda^{(\omega+1)/2})^2$ would be a square, a contradiction; hence $\omega$ is even, and so $T=\alpha\omega$ is even.

The final assertion follows: on a full period, \cref{cor:quadratic-half-period} requires $T$ even and $\chiP(-1)=-1$, that is, $p\equiv3\pmod4$, and then $\chiP(\Lambda)=-1$ by (b), so \cref{cor:quadratic-scalar} applies.
\end{proof}

\begin{corollary}[Quadratic full-period configurations]
In the inert setting, for quadratic-character sums over a full matrix period, the half-period quadratic cancellation criterion implies the rank-block scalar cancellation criterion. The reverse implication need not hold. Consequently, among the four formal configurations of the two criteria, the configuration in which half-period quadratic cancellation holds but scalar cancellation fails is impossible.
\end{corollary}

\begin{proof}
For a full period, the half-period quadratic criterion requires $T$ even and $\chi_p(-1)=-1$, equivalently $p\equiv 3\pmod 4$. Proposition 5.7 then gives $\chi_p(\Lambda)=-1$, so Corollary 5.6 applies. The failure of the reverse implication is shown by Example 6.2, where $\chi_p(\Lambda)=-1$ but $\chi_p(-1)=1$.
\end{proof}


\begin{remark}[Two mechanisms and their relation]
\label{rem:two-mechanisms}
The cancellation in \cref{cor:quadratic-half-period} is a translation pairing:
\[
 X_{n+T/2}=-X_n.
\]
It requires an appropriate index-set symmetry and an odd character, or more generally an odd test function. By contrast, the cancellation in \cref{cor:quadratic-scalar} is rank-block scalar propagation, equivalently a geometric character factorization:
\[
 X_{k\alpha+j}=\Lambda^kX_j.
\]
It requires $\psi(\Lambda)\neq1$ and does not require $\psi(-1)=-1$. Thus scalar cancellation can occur even when the half-period negation pairing is ineffective, as in \cref{ex:scalar-cancellation}.

For the quadratic character on a full matrix period, \cref{prop:chi-lambda} nevertheless shows that the two criteria are related: the half-period criterion implies the scalar criterion, while the converse need not hold. Moreover, the scalar criterion also forces the existence of a half-period involution:
indeed, $\chiP(\Lambda)=-1$ implies that $\omega$ is even and hence
$T=\alpha\omega$ is even, so that $A^{T/2}=-I$ by
\cref{prop:half-period-criterion}. However, the resulting anti-periodicity
produces quadratic-character cancellation only when
$\chiP(-1)=-1$.
\end{remark}

\section{Examples}

\begin{example}[Fibonacci sequence modulo $7$]
\label{ex:fibonacci-7}
Take
\[
 (P,Q,p)=(1,-1,7).
\]
Then
\[
 D=P^2-4Q=5,
\]
which is a nonsquare modulo $7$. Hence $X^2-X-1$ is irreducible over $\mathbb{F}_7$. The companion matrix is
\[
 A=
 \begin{pmatrix}
 1&1\\
 1&0
 \end{pmatrix}.
\]
Direct exponentiation gives
\[
 A^8=-I,\qquad A^{16}=I,
\]
so
\[
 T=16,\qquad U_{n+8}=-U_n\pmod7.
\]
The rank data provide an independent cross-check:
\[
 \alpha=8=p+1,\qquad
 \Lambda=U_9\equiv-1\pmod7,\qquad
 \omega=\ord_{\mathbb{F}_7^\times}(-1)=2,\qquad
 T=\alpha\omega=16.
\]
Here $\chiP(-1)=\chiP(\Lambda)=-1$, so both mechanisms are active: the half-period translation gives
\[
 \chiP(U_{n+8})=-\chiP(U_n),
\]
as in \cref{cor:quadratic-half-period}, and the same vanishing also follows from \cref{cor:quadratic-scalar}.
\end{example}

\begin{example}[Inert scalar cancellation]
\label{ex:scalar-cancellation}
Take
\[
 (P,Q,p)=(1,2,5).
\]
Here
\[
 D=1-8=-7\equiv3\pmod5,
\]
and $3$ is a nonsquare modulo $5$. Hence $X^2-X+2$ is irreducible over $\mathbb{F}_5$. The initial values are
\[
 U_0,U_1,U_2,U_3,U_4,U_5,U_6,U_7
 =
 0,1,1,4,2,4,0,2
 \pmod5.
\]
The rank and scalar multiplier are
\[
 \alpha=6,\qquad \Lambda=U_{\alpha+1}=U_7=2.
\]
Moreover,
\[
 \chiP(\Lambda)=\chi_5(2)=-1,
 \qquad
 \omega=\ord_{\mathbb{F}_5^\times}(2)=4,
 \qquad
 T=\alpha\omega=24.
\]
Since $\alpha=6$ is even, the relation $\chiP(\Lambda)=-1$ is consistent with \cref{prop:chi-lambda}(a). The identity $A^\alpha=\Lambda I$ reads $A^6=2I$, whence $A^{12}=-I$: anti-periodicity is present. However,
\[
 \chiP(-1)=\chi_5(-1)=1,
\]
so the half-period translation does not itself produce quadratic-character cancellation. Instead, \cref{prop:rank-multiplier} gives
\[
 U_{6k+j}=2^kU_j\pmod5,
 \qquad
 V_{6k+j}=2^kV_j\pmod5.
\]
The four rank blocks have quadratic-character sums
\[
 3,\ -3,\ 3,\ -3,
\]
and the alternating geometric factor $\chiP(\Lambda)^k=(-1)^k$ makes the total vanish. This is scalar character filtering, not a negation pairing.
\end{example}

\begin{remark}[Even period alone is not enough]
\label{rem:even-period-not-enough}
The conditions $T$ even and $A^{T/2}=-I$ alone do not force quadratic-character cancellation. For $(P,Q,p)=(1,1,5)$ one has $D\equiv2\pmod5$ (a nonsquare), rank $\alpha=3$, multiplier $\Lambda\equiv-1\equiv4\pmod5$, order $\omega=2$, and period $T=6$; then $A^3=-I$, but $\chi_5(-1)=\chi_5(\Lambda)=1$ and
\[
 \sum_{n=1}^{6}\chi_5(U_n)=4\neq0.
\]
Together, \cref{ex:fibonacci-7}, \cref{ex:scalar-cancellation}, and this remark realize the three configurations of the two quadratic criteria on a full period that can occur: both active, scalar only, and neither. By \cref{prop:chi-lambda}, the fourth configuration---half-period cancellation without scalar cancellation---cannot occur.
\end{remark}

\begin{table}[ht]
\centering
\small
\begin{tabular}{c c c c c c c c}
\hline
Case & $(P,Q,p)$ & $T$ even? & $\chi_p(-1)$ & $\chi_p(\Lambda)$ & $HP$ quadratic & Scalar & Active criteria \\
\hline
Example 6.1 & $(1,-1,7)$ & yes & $-1$ & $-1$ & yes & yes & both \\
Example 6.2 & $(1,2,5)$ & yes & $1$ & $-1$ & no & yes & scalar only \\
Remark 6.3  & $(1,1,5)$ & yes & $1$ & $1$ & no & no & neither \\
\hline
\end{tabular}
\caption{The three attainable configurations represented in the examples. The fourth formal configuration, half-period quadratic cancellation without scalar cancellation, is excluded by the quadratic full-period configuration result above.}
\end{table}


\section{Scope Remark}

The purpose of this note is to isolate cancellation mechanisms rather than to modify the classical rank--period theory. The rank--period structure supplies the algebraic scaffolding, while the contribution here is the explicit separation of translation cancellation from scalar character filtering. The decomposition into rank blocks follows from the companion-matrix identity, the definition of the rank of apparition, and the divisibility $\alpha\mid n$ whenever $U_n=0\pmod p$, which in the inert setting is an immediate consequence of \cref{lem:rank-characterization}. For the quadratic character on a full period these criteria are related as in \cref{prop:chi-lambda}, but they remain distinct as proofs.

The hypotheses are essential. If $p\mid Q$, the companion matrix need not be invertible. If $p\mid D$, the characteristic polynomial is not separable in the required form. In either situation, the inert quadratic-extension argument used here does not apply without modification. Outside the inert setting considered here, an even matrix period alone does not in general imply that the half-period matrix is $-I$; within the inert setting, $T$ is even if and only if $A^{T/2}=-I$, by \cref{prop:half-period-criterion}.

\section*{Acknowledgements}

The author acknowledges the use of artificial intelligence (AI) tools for
linguistic editing, organizational refinement, notation consistency, and
LaTeX assistance during the preparation of the manuscript. The author independently developed, verified, and assumes full responsibility for all
mathematical statements, proofs, examples, citations, and computational
checks presented in the manuscript.


\section*{Funding}

No funding was received for this work.

\section*{Disclosure}

The author declares no conflict of interest.
\section*{Data Availability}

No external datasets were used. The numerical checks and examples are
reproducible directly from the displayed recurrences, companion matrices,
and finite-field computations. The verification code is archived in the
Zenodo reproducibility release \cite{Ghandali2026HalfPeriodSoftware}. The
archive contains the pure-Python scripts
\texttt{code/verify-examples.py} and
\texttt{code/verify-theorem-pipeline.py}, which verify, for the parameter
triples displayed in Section~6, the rank of apparition, the multiplier, the
matrix period, the half-period identity, and the values of the stated
character sums, including their vanishing in
\cref{ex:fibonacci-7,ex:scalar-cancellation} and their non-vanishing in
\cref{rem:even-period-not-enough}.

A public preprint version of this manuscript is archived at
\url{https://doi.org/10.5281/zenodo.22149586}.
 

\bibliographystyle{plain}
\bibliography{references}

\end{document}
